
\de{ĐỀ THI GIỮA HỌC KỲ II NĂM HỌC 2022-2023}{THPT Bình Phú}

\begin{bt}%[0T7Y1-1]%[0T7B2-1]%[0T7K2-2]%[Dự án đề kiểm tra HKII NH22-23-Thy Nguyen Vo Diem]%[Bình Phú]
\begin{enumerate}
	\item Lập bảng xét dấu tam thức $f(x)=-x^2+5 x-6$.
	\item Cho $g(x)=m x^2-2 m x-2$. Xác định $m$ để $g(x) \leq 0$,  $\forall x \in  \mathbb{R}$.
	\item Cho $h(x)=x^2\left(9-x^2\right)$. Xác định số nguyên $a$ lớn nhất thỏa $h(a)>0$.
\end{enumerate}
\loigiai{
\begin{enumerate}
	\item Cho $f(x)=0 \Leftrightarrow \hoac{&x=3\\&x=2}$. Do $a=-1<0$ nên ta có bảng xét dấu $f(x)$
	\begin{center}
\begin{tikzpicture}
\tkzTabInit[lgt=1.2,espcl=3]
{$x$ /0.7, $f(x)$ /0.7}
{$-\infty$,$2$,$3$,$+\infty$}
\tkzTabLine{ ,-,0,+,0,-, }
\end{tikzpicture}
\end{center}
	\item \textbf{TH1:} $m=0$, ta có $g(x)=-2 \le 0$, $\forall x \in \mathbb{R}$. Vậy $m=0$ thỏa đề.\\
	\textbf{TH2:} $m \ne 0$.
	\allowdisplaybreaks{\begin{eqnarray*}
		g(x) \leq 0, \, \forall x \in \mathbb{R}	&\Leftrightarrow& m x^2-2 m x-2 \leq 0, \, \forall x \in \mathbb{R} \\
			&\Leftrightarrow& \heva{&\Delta'\le0\\&a<0} \\
			&\Leftrightarrow& \heva{&m^2+2m\le0\\&m<0} \\
			&\Leftrightarrow& \heva{&-2 \le m \le 0\\&m<0}\\
			&\Leftrightarrow& -2 \le m < 0.
	\end{eqnarray*}}
Kết hợp điều kiện $m \ne 0$, ta có $-2 \le m < 0$.\\
Vậy $-2 \le m \le 0$ thỏa yêu cầu bài toán.
	\item Xét $h(x)=x^2\left(9-x^2\right)$ có $x^2 \ge 0$, $\forall x \in \mathbb{R}$ nên 
	\[h(x)>0	\Leftrightarrow \heva{&x \ne 0\\& 9-x^2>0} \Leftrightarrow \heva{&x \ne 0\\& -3<x<3} \Leftrightarrow \hoac{&-3<x<0\\&0<x<3}.\]
	Vậy số nguyên $a$ lớn nhất thỏa $h(a)>0$ là $a=2$.
\end{enumerate}
}
\end{bt}

\begin{bt}%[0T7Y2-1]%[0T7B2-2]%[0T7B3-2]%[0T7K3-2]%[Dự án đề kiểm tra HKII NH22-23-Thy Nguyen Vo Diem]%[Bình Phú]
	Giải các phương trình, bất phương trình sau
	\begin{listEX}[2]
	\item $3x^2-5x+2 \geq 0$
	\item $\dfrac{-x^2}{x^2-4}+1<0$
	\item $\sqrt{2x+1}=\sqrt{x^2+2 x+1}$
	\item $(x-3)\left(\sqrt{8-x^2}-x\right)=0$.
	\end{listEX}
	\loigiai{
		\begin{listEX}[1]
			\item $3x^2-5x+2 \geq 0 \Leftrightarrow \hoac{&x\le \dfrac{2}{3}\\&x \ge 1}$.\\
			Vậy tập nghiệm $S=\left(-\infty;\dfrac{2}{3}\right]\cup [1;+\infty)$.
			\item $\dfrac{-x^2}{x^2-4}+1<0 \Leftrightarrow \dfrac{-4}{x^2-4}<0 \Leftrightarrow x^2-4>0 \Leftrightarrow \hoac{&x>2\\&x<-2}$.\\
			Vậy tập nghiệm $S=\left(-\infty;-2\right)\cup (2;+\infty)$.
			\item $\sqrt{2x+1}=\sqrt{x^2+2 x+1} \Leftrightarrow \heva{&2x+1 \ge 0\\&2x+1=x^2+2 x+1} \Leftrightarrow \heva{&x \ge -\dfrac{1}{2}\\&x^2=0} \Leftrightarrow x=0$.\\
			Vậy tập nghiệm $S=\{0\}$.
			\item Ta có \allowdisplaybreaks{\begin{eqnarray*}
				(x-3)\left(\sqrt{8-x^2}-x\right)=0	&\Leftrightarrow&  \hoac{&x-3=0\\&\sqrt{8-x^2}-x=0}\\
					&\Leftrightarrow& \hoac{&x=3\\&\sqrt{8-x^2}=x} \\
					&\Leftrightarrow& \hoac{&x=3\\&\heva{&x\ge 0\\&8-x^2=x^2}} \\
					&\Leftrightarrow& \hoac{&x=3\\&\heva{&x\ge 0\\&x=\pm 2}} \\
					&\Leftrightarrow& \hoac{&x=3\\&x=2.} 
			\end{eqnarray*}}
		\end{listEX}
	Vậy tập nghiệm $S=\{2;3\}$.
	}
\end{bt}


\begin{bt}%[0T7T2-1]%[Dự án đề kiểm tra GHKI NH22-23 - Lê Hùng Thắng]%%[THPT Bình Phú - HCM]%Bài 3
	Chị Hạnh bán trà sữa tại nhà. Chị dự tính tổng chi phí $M$ để pha chế $n$ ly trà sữa được cho bởi biểu thức $M=n^2-110 n+3300$ (nghìn đồng). Giả sử tất cả các ly trà sữa được pha chế thì bán hết với giá bán của một ly là $30000$ đồng. Hỏi số ly trà sữa được pha chế phải đạt trong khoảng nào để chị Hạnh không bị lỗ?
	\loigiai{
	Theo đề ta có\\
	Tổng chi phí cho $n$ ly trà sữa là $M=n^2-110 n+3300$ (nghìn).\\
	Tổng số tiền khi bán hết $n$ ly trà sữa là $T=30 n$ (nghìn).\\
	Việc bán trà sữa không bị lỗ vốn khi
	\allowdisplaybreaks \vspace*{-0.4cm}
	$$\begin{aligned}[t]
	M \le T \Leftrightarrow &\, n^2-110 n+3300 \le 30 n \\
	 \Leftrightarrow &\, n^2- 140 n +3300 \le 0 \\
	 \Leftrightarrow &\, 30 \le n \le 70. 
	\end{aligned}$$
Vậy để không bị lỗ thì số ly trà sữa có thể pha chế để bán là $n$ thỏa $30 \le n \le 70$. 
	}
\end{bt}

%D:\LaTeX Nhien Hoai\Du an lam de LaTeX lht\0-TL-BinhPhu-GHKI-NH22-23-cau3-4.tex
\begin{bt}%[0T9Y1-3]%[0T9B2-2]%[0T9B3-2]%[0T9G2-2]%[Dự án đề kiểm tra GHKI NH22-23 - Lê Hùng Thắng]%%[THPT Bình Phú - HCM]%Bài 4
Trong mặt phẳng tọa độ $Oxy$, cho $\triangle ABC$ có $A(0;4), B(-4;2), C(1;2)$.
\begin{enumerate}
	\item Xác định tọa độ trung điểm $I$ của $AB$ và trọng tâm $G$ của $\triangle ABC$.
	\item Viết phương trình tổng quát của đường thẳng $AB$.
	\item Viết phương trình đường tròn $(C)$ có tâm $A$ và đi qua $C$.
	\item Viết phương trình đường thẳng $d$ qua điểm $C$ đồng thời cắt $Ox$ và $Oy$ lần lượt tại hai điểm $E$ và $F$ sao cho $\triangle OEF$ có diện tích bằng $\dfrac{1}{2}$.
\end{enumerate}
	\loigiai{
	\begin{enumerate}
	\item Trung điểm của $AB$ là $I(-2;3)$.\\
	Trọng tâm của $\triangle ABC$ là $G\left(-1;\dfrac{8}{3}\right)$.
	\item Đường thẳng $AB$ có véc-tơ chỉ phương là $\vec{u}=\vec{AB}=(-4;-2)$.\\
	$AB$ đi qua $A$ và có véc-tơ pháp tuyến $\vec{n}=(1;-2)$ nên $AB$ có phương trình là
	$$1(x-0)-2(y-4)=0 \Leftrightarrow x-2y+8=0.$$
	\item Vì $(C)$ có tâm $A$ và đi qua $C$ nên $(C)$ có bán kính là $R=AC=\sqrt{(1-0)^2+(2-4)^2}=\sqrt{5}$.\\
	Vậy $(C)$ có phương trình là $x^2+(y-4)^2=5$.
	\immini{
	\item Gọi $E(a;0) \in Ox$ và $F(0;b)\in Oy$.\\
	Vì $d$ qua $E$ và $F$ nên phương trình của $d$ theo đoạn chắn là $\dfrac{x}{a}+\dfrac{y}{b}=1$.\\
	Mà $d$ qua $C(1;2)$ nên ta có $\dfrac{1}{a}+\dfrac{2}{b}=1 \Leftrightarrow 2a+b=a\cdot b$ \qquad (1).\\
	\allowdisplaybreaks
	$\begin{aligned}[t]
	 S_{\triangle OEF}= \dfrac{1}{2}
	 \Leftrightarrow & \dfrac{1}{2} OE \cdot OF = \dfrac{1}{2}\\
 \Leftrightarrow & \dfrac{1}{2} |a|\cdot |b| =\dfrac{1}{2}\\
 \Leftrightarrow & |a\cdot b| =1 \Leftrightarrow \hoac{&a\cdot b= 1 \,\,\quad (2) \\& a\cdot b= -1 \quad (3).}
	\end{aligned}$
	}
	{\begin{tikzpicture}[scale=1,>=stealth, font=\footnotesize, line join=round, line cap=round]
	\def\xmin{-2} \def\xmax{2}
	\def\ymin{-1} \def\ymax{3}
%	\draw[color=gray!30,dashed] (\xmin,\ymin) grid (\xmax,\ymax);
	\draw[->] (\xmin,0)--(\xmax,0) node [below]{$x$};
	\draw[->] (0,\ymin)--(0,\ymax) node [right]{$y$};
	\node at (0,0) [below right]{$O$};
	\clip (\xmin+0.1,\ymin+0.1) rectangle (\xmax-0.1,\ymax-0.1);
	\draw[smooth,samples=200] plot(\x,{1*(\x)+1});
	\draw [dotted] (1,0)|-(0,2);
	\fill[pattern=north east lines,opacity=0.5](0,0)--(-1,0)--(0,1)--cycle;
	\fill (1,0) circle (1pt) node [below] {$1$}
	(1,2) circle (1pt) node [right] {$C$}
	(-1,0) circle (1pt) node [below] {$E$}
	(0,1) circle (1pt) node [left] {$F$}
	(0,2) circle (1pt) node [left] {$2$};
	\end{tikzpicture}}
\noindent 
Từ (1) và (2) ta có $\heva{&2a+b=1 \\&a\cdot b=1}  \Leftrightarrow \heva{&b=1-2a \\& a(1-2a)=1} \Leftrightarrow \heva{&b=1-2a \\&-2a^2+a-1=0.}$ (hệ vô nghiệm).\\
Từ (1) và (3) ta có $\heva{&2a+b=-1 \\&a\cdot b=-1}  \Leftrightarrow \heva{&b=-2a-1 \\& a(-2a-1)=-1} \Leftrightarrow \heva{&b=-2a-1 \\&2a^2+a-1=0} \Leftrightarrow \heva{& b=-2a-1 \\& \hoac{&a=-1 \\& a= \dfrac{1}{2}.}}$\\
$\bullet$ Với $a=-1 \Rightarrow b=1$ khi đó $d\colon x-y+1=0$.\\
$\bullet$ Với $a=\dfrac{1}{2} \Rightarrow b=-2$ khi đó $d\colon 4x-y-2=0$.\\
Vậy có hai đường thẳng $d$ thỏa đề bài.
	\end{enumerate}	
	}
\end{bt}


